\chapter{שיטת עבודה ומתודולוגיה}\label{ch:method}

\section{מודל הפיתוח}
הפרויקט משלב שני מודלים: מודל \en{V} ברמת המערכת, לפי \en{ISO/IEC/IEEE 15288}~\cite{iso15288}, וספרינטים של שבועיים ברמת הרכיב. מודל \en{V} נבחר משום שהמוצר חייב להוכיח עמידה בדרישות מוגדרות מראש: כל שלב תכנון בצד שמאל של ה-\en{V} מקבל שלב אימות מקביל בצד ימין. הספרינטים נבחרו כי גלאים ומודלי כיול דורשים ניסוי וטעייה מהירים (איור~\ref{fig:vmodel}).

\begin{figure}[!htbp]
\begin{ltrdiagram}
% vmodel: TikZ / pgfplots source, built inside ltrdiagram (LTR)
\begin{tikzpicture}[font=\sffamily\scriptsize,
  bx/.style={draw=acc,fill=soft,rounded corners=3pt,minimum width=3.1cm,minimum height=0.75cm,align=center},
  vb/.style={draw=okc,fill=okc!10,rounded corners=3pt,minimum width=3.1cm,minimum height=0.75cm,align=center},
  ar/.style={-{Stealth[length=5pt]},thick,ink!60}]
\node[bx] (l1) at (0,3.6) {\he{דרישות לקוח ותקן אתי}};
\node[bx] (l2) at (1.3,2.4) {\he{דרישות מערכת, פרק 6}};
\node[bx] (l3) at (2.6,1.2) {\he{ארכיטקטורה, פרק 8}};
\node[draw=acc2,fill=soft2,rounded corners=3pt,minimum width=3.1cm,minimum height=0.75cm,align=center] (b) at (5.2,0) {\he{מימוש בספרינטים}};
\node[vb] (r3) at (7.8,1.2) {\he{בדיקות שילוב}};
\node[vb] (r2) at (9.1,2.4) {\he{בדיקות מערכת על \en{VL-Wild}}};
\node[vb] (r1) at (10.4,3.6) {\he{פיילוט וקבלה}};
\draw[ar] (l1)--(l2); \draw[ar] (l2)--(l3); \draw[ar] (l3)--(b);
\draw[ar] (b)--(r3); \draw[ar] (r3)--(r2); \draw[ar] (r2)--(r1);
\draw[dashed,gray,<->] (l1)--(r1) node[midway,above]{\he{קריטריוני קבלה}};
\draw[dashed,gray,<->] (l2)--(r2) node[midway,above]{\he{מטריצת עקיבות}};
\draw[dashed,gray,<->] (l3)--(r3) node[midway,above]{\he{חוזי ממשק}};
\end{tikzpicture}

\end{ltrdiagram}
\caption{מודל \en{V}: לכל שלב תכנון יש שלב בדיקה מקביל שמוגדר כבר עכשיו}\label{fig:vmodel}
\end{figure}

\section{מתודולוגיית המחקר והפיתוח של הגלאים}
\begin{enumerate}
\item \textbf{איסוף נתונים:} מאגר אימון של \num{1.2} מיליון תמונות: מחצית צילומים אמיתיים מסוכנויות שותפות ברישיון, ומחצית תמונות סינתטיות מ-\num{12} מחוללים מסחריים ופתוחים. כל תמונה סינתטית נוצרת מכיתוב של צילום אמיתי, כדי שהגלאי ילמד את ההבדל בין אמת לזיוף ולא את ההבדל בין נושאים.
\item \textbf{הדמיית ״השטח״:} כל תמונה עוברת אקראית שרשרת עיבוד שמחקה הפצה ברשת: דחיסה חוזרת, הקטנה, צילום מסך, חיתוך והוספת טקסט. זה המענה הישיר לממצא שגלאים קורסים מחוץ למעבדה~\cite{ren2025}.
\item \textbf{הפרדה לפי מחולל:} שלושה מחוללים מוצאים מהאימון לחלוטין ומשמשים רק לבדיקה. כך נמדד הביצוע מול מחולל ״חדש״ שהמערכת לא ראתה.
\item \textbf{כיול:} מאגר נפרד של \num{50,000} תמונות משמש רק לכיול הציון.
\item \textbf{אימון חוזר:} בכל רבעון מתווספים מחוללים חדשים. גרסה חדשה מוחלפת רק אם אינה גרועה מקודמתה בשום מדד.
\end{enumerate}

\section{ועדת המומחים}\label{sec:committee}
לפני הכתיבה הגדרנו עם סוכן בינה מלאכותית ״ועדת מומחים״ של שבעה תפקידים: מנהל פרויקטים ומעריך מכרזים (יו״ר), מהנדס מערכות, מומחה תחום, מנהל סיכונים ואיכות, אנליסט עסקי, חוקר אקדמי ועורך טכני. כל תפקיד כתב את הדרישות לפרקים שבאחריותו, ביקר את דרישות האחרים, והיו״ר ניקד את הטיוטה לפי מחוון. התקן שהוועדה כתבה הוא הבסיס למבנה הדוח (הפרומפט המלא בתיקייה \code{prompts/}).

\section{מתודולוגיית פיתוח התקן האתי}
התקן נכתב בתהליך של ארבעה שלבים: (\num{1}) ניתוח של \num{40} מקרי מבחן מתועדים לפי סוג הנזק; (\num{2}) טיוטה ראשונה של ועדת המומחים; (\num{3}) שלוש סדנאות עם עורכים ממערכות הפיילוט, שבהן כל כלל נבחן מול מקרה אמיתי מהמשמרת; (\num{4}) סבב הערות ציבורי של \num{30} יום. כלל שעורכים לא הצליחו ליישם בסדנה נכתב מחדש או נמחק.

\section{כלים וסביבת עבודה}

ראו טבלה~\ref{tab:tools}.

\begin{table}[!htbp]\centering
\caption{כלים וסביבות}\label{tab:tools}
\begin{tabularx}{\linewidth}{@{}l Y@{}}
\toprule
\textbf{תחום} & \textbf{כלי} \\
\midrule
קוד ובקרת גרסאות & \en{GitHub}, \en{CI/CD} עם בדיקות אוטומטיות בכל מיזוג \\
למידת מכונה & \en{PyTorch}, \en{Weights \& Biases} למעקב ניסויים \\
תשתית & \en{Kubernetes} בענן, \en{Terraform} \\
ניהול פרויקט & \en{Jira} (ספרינטים), \en{Confluence} (תיעוד), גאנט ב-\en{MS Project} \\
מסמכים & \en{LuaLaTeX} עם תבנית משלנו (\code{my-template.cls}), מקורות ב-\en{biblatex} ו-\en{biber}, ארבע הרצות הידור \\
שימוש בבינה מלאכותית בכתיבת הדוח & סוכן בינה מלאכותית בטרמינל כתב את קובצי הלטך לפי הנחיות שלנו; הפרומפטים שמורים וממוספרים בתיקייה \code{prompts/}, וכל נתון ומקור נבדק מול דף המקור \\
\bottomrule
\end{tabularx}
\end{table}

\section{איך הדוח עצמו נבנה ונבדק}\label{sec:qa}
הדוח נבנה כפרויקט ולא כקובץ אחד, לפי עקרונות הקורס~\cite{segal2026l03}:
\begin{itemize}
\item \textbf{מבנה:} קובץ ראשי (\code{main.tex}) שמחבר את הפרקים, פרק בכל קובץ (\code{chapters/}), כל תרשים בקובץ נפרד (\code{figures/}), רשימת מקורות כבסיס נתונים (\code{references.bib}), וטיוטות ופרומפטים מחוץ למסמך (\code{work/}, \code{prompts/}).
\item \textbf{תבנית:} כל העיצוב נמצא ב-\code{my-template.cls} (גרסה \num{1.0}, עם יומן שינויים ורשימת פקודות ב-\code{TEMPLATE.md}). בתוך הפרקים אין עיצוב, רק פקודות של התבנית.
\item \textbf{עברית ואנגלית:} כל מילה באנגלית מסומנת ב-\code{\textbackslash en}, כל מספר, אחוז וגרסה ב-\code{\textbackslash num}, וכל תרשים נבנה בסביבה אנגלית.
\item \textbf{בקרת איכות:} אחרי ההידור עבר הדוח סבב של סוכני בדיקה נפרדים, שכל אחד מהם בדק משהו אחר. הממצאים והתיקונים מתועדים ב-\code{work/review/}, ומסוכמים בטבלה~\ref{tab:review}.
\end{itemize}

% review-summary.tex — סיכום סבב סוכני הבדיקה. הדוחות המלאים: work/review/
טבלה~\ref{tab:review} מסכמת את ארבעת סבבי הבדיקה. כל סוכן קיבל פרומפט נפרד ובדק משהו אחר, והסוכן שכתב את הדוח לא בדק את עצמו. הסוכנים רק דיווחו; את התיקונים עשינו אחר כך, וקימפלנו מחדש.

\begin{table}[H]\centering
\caption{סוכני הבדיקה: מה כל אחד בדק, מה מצא ומה תוקן}\label{tab:review}
\small
\begin{tabularx}{\linewidth}{@{}l Y l Y@{}}
\toprule
\colhead{סוכן} & \colhead{מה בדק} & \colhead{ממצאים} & \colhead{מה תוקן} \\
\midrule
כיווניות & עברית ואנגלית, עמוד אחר עמוד: אותיות הפוכות, מספרים, סוגריים, תוכן עניינים & \num{22} & סימן כיוון בלתי נראה בפקודות התבנית (גרסה \num{1.1}); בלי סוגריים בתוויות של תרשימים \\
עמוד אחר עמוד & כותרות בתחתית עמוד, עמודים ריקים, תרשימים חופפים, טבלאות ואיורים בלי הפניה & \num{24} & הפניה לכל טבלה ואיור; תוויות וגרפים שתוקנו; כותרת לא נשארת לבד בתחתית עמוד \\
מקורות & כל מקור קיים, והמספר בטקסט באמת מופיע במקור & \num{15} & נתוני \en{NewsGuard} ו-\en{Nikon} תוקנו; מבחן ארבעת הרכיבים יוחס לטיוטת ההנחיות; פרטי ביבליוגרפיה תוקנו \\
צוות אדום & ציון לפי דרישות המטלה, עקביות המספרים, היקף החוזה & ציון \num{56} מתוך \num{100} & \num{8} סתירות מספריות תוקנו; נוסף ״היקף: חזון מול מטלה״; הוסרו מספרי תעודת זהות מהמסמך \\
בדיקה חוזרת & אחרי התיקונים: האם כל ממצא תוקן, ומה חדש & נותרו \num{9} & כל הממצאים הקריטיים תוקנו; עמודים כמעט ריקים, מספור טבלאות בנספחים ופיצול מילה תוקנו בסבב נוסף \\
ציון סופי & גרסה \num{2.0} כולה: דרישות המטלה, מאגר ותהליך, קול אישי, מקורות & ציון \num{68} מתוך \num{100} & מקורות ראשוניים במקום בלוגים ו-\num{4} מאמרים אקדמיים; סתירת ״הפער״ ולשון עבר בסעיף~\ref{sec:ethicsmethod} תוקנו; מספרי תעודת זהות נמחקו גם מהיסטוריית המאגר; התבנית בגרסה \num{1.2} \\
ציון חוזר & גרסה \num{2.1} כולה, כולל סעיפים \ref{sec:enforce} ו-\ref{sec:israel}; \num{7} נתונים מול המקורות & ציון \num{76} מתוך \num{100} & כל \num{7} הנתונים נמצאו נכונים; איורים, גרסת התבנית, גרסת המסמך וצילומי המסך תוקנו. נותרו: הסעיף האישי של פז וה-\en{commits} שלו \\
הידור (בדיקה שלנו) & שגיאות, הפניות חסרות ושורות חורגות ביומן ההידור & \num{3} & הפניה חסרה נוספה; שורות ארוכות נשברות \\
\bottomrule
\end{tabularx}
\end{table}

